% 总体模型、随机样本、统计量（估计量）与估计值的关系示意。
% 上行（蓝）是抽样前的随机对象，下行（绿）是一次抽样后的数值；纵向箭头表示取一次实现。
\begin{tikzpicture}[line width=.6pt,font=\fontsize{8}{10}\selectfont,>=Stealth,
  obj/.style={draw=FigureBlue,fill=FigureBlueFill,rounded corners=1.8pt,align=center,
    minimum width=26mm,minimum height=11mm,inner sep=2pt,text=FigureInk},
  val/.style={draw=FigureGreen,fill=FigureGreenFill,rounded corners=1.8pt,align=center,
    minimum width=26mm,minimum height=11mm,inner sep=2pt,text=FigureInk},
  pop/.style={draw=FigureYellow,fill=FigureYellowFill,rounded corners=1.8pt,align=center,
    minimum width=26mm,minimum height=11mm,inner sep=2pt,text=FigureInk},
  arr/.style={->,draw=FigureStroke,line width=.6pt},
  lab/.style={font=\fontsize{7.5}{9}\selectfont,text=FigureMuted}]
\node[pop] (P) at (0,0) {总体模型\\$F_\theta,\ \theta\in\Theta$};
\node[obj] (S) at (41mm,0) {随机样本\\$X_1,\dots,X_n$};
\node[obj] (T) at (82mm,0) {统计量（估计量）\\$T=T(X_1,\dots,X_n)$};
\node[val] (x) at (41mm,-21mm) {观测数据\\$x_1,\dots,x_n$};
\node[val] (t) at (82mm,-21mm) {估计值\\$t=T(x_1,\dots,x_n)$};
\draw[arr] (P) -- node[above,lab]{独立抽取} (S);
\draw[arr] (S) -- node[above,lab]{规则$T$} (T);
\draw[arr] (x) -- node[above,lab]{同一规则$T$} (t);
\draw[arr] (S) -- node[left,lab]{一次实现} (x);
\draw[arr] (T) -- node[right,lab]{代入数据} (t);
\draw[->,draw=FigureRed,dashed,line width=.6pt] (t.south) -- ++(0,-5mm) -| node[pos=.25,below,lab]{用估计值猜测未知参数$\theta$} (P.south);
\end{tikzpicture}
